\par\smallskip\noindent\textbf{Abbreviations used in this section:} ACL --- Access control list; API --- Application programming interface; CSPRNG --- Cryptographically secure pseudorandom number generator; HMAC --- Hash-based message authentication code; ID --- Identifier; JSON --- JavaScript Object Notation; JWT --- JSON Web Token; SHA-256 --- Secure Hash Algorithm with a 256-bit digest; TLS --- Transport Layer Security; TTL --- Time to live; YAML --- YAML Ain't Markup Language (recursive name).\par\smallskip
% API credentials on Redis: included from catalyst-sessions.tex
\section{API credentials on Redis}\label{sec:api-tokens}
Browsers authenticate with a session cookie. API clients, such as mobile apps, single-page applications, partner integrations and command-line tools, send a credential in a header with every request. The OpenAPI document describes those credentials as \emph{security schemes}; Redis is where the server checks, rotates, limits and revokes them. Every concurrency problem in this handbook turns up here too: two tabs refreshing the same token are a race, a retried signed request is a replay, and a per-key quota is a shared counter.

Table~\ref{tab:credentials} compares the five credentials the laboratory implements. They complement one another: a browser logs in with a session, creates an API key for a script, or obtains a short-lived access token and a refresh token for a mobile client.

\begin{table}[!htbp]\centering\small
\caption{Credentials, their lifetimes and where Redis takes part.}\label{tab:credentials}
\begin{tabular}{L{.14\linewidth}L{.14\linewidth}L{.12\linewidth}L{.27\linewidth}L{.17\linewidth}}\toprule
Credential & Held by & Lifetime & Redis keys and checks & OpenAPI scheme \\ \midrule
Session cookie & browser & idle timeout & \code{sess:<id>}, per-session lock & \code{apiKey}, \code{in: cookie} \\
API key & scripts, integrations & months; rotated & \code{apikey:<sha256>}; per-key rate counter & \code{apiKey}, \code{in: header} \\
Access token (JWT) & API client & 5 minutes & none on the happy path; \code{jwt:deny:<jti>}, family check & \code{http}, \code{bearer}, \code{JWT} \\
Refresh token & API client, stored securely & 30 days; rotated on every use & \code{rt:{fam}:<sha256>}, \code{rtfam:{fam}} & issued by \code{/token}; \code{oauth2} if an authorization server issues it \\
Signed request & partner systems & one request & \code{nonce:<kid>:<nonce>} (\code{SET NX}) & \code{apiKey} header, described \\
\bottomrule\end{tabular}
\end{table}

\subsection*{Describing the schemes in OpenAPI 3.2}
The laboratory's \code{openapi.yaml} declares four schemes and states the requirement on every operation. An empty requirement list (\code{security: []}) marks an operation as public, which is clearer than relying on a global default. Alternatives in a requirement list are ORed; schemes inside one requirement object are ANDed.
\CodeFile{yaml}{snippets/api-credentials-01.yaml}{Security schemes and one operation from \texttt{openapi.yaml}}
\version{OpenAPI 3.2.0 (September 2025) [51] added three security features: the OAuth~2.0 device authorization flow (\code{deviceAuthorization}), \code{oauth2MetadataUrl} for discovering an authorization server's metadata, and a \code{deprecated} flag that tells clients to move off an old scheme, such as a legacy API-key header during a migration to tokens. An OpenAPI document describes credentials; it does not enforce them. \code{OpenAPI::Modern} [53] can validate requests and responses against the document in tests and middleware. The laboratory validates its document with \code{openapi-spec-validator} and checks that every exercised route is documented with an explicit requirement; it does not use \code{OpenAPI::Modern} at run time.}
\tip{Put the OpenAPI document under test. Example 30 loads it with \code{YAML::XS} and fails if a route it calls is missing or has no \code{security} entry. A forgotten requirement then shows up as a failing test, not as an unprotected endpoint.}

% ---------------------------------------------------------------- API keys
\subsection*{API keys: store the hash, never the key}
\index{API!sha256_hex@\texttt{sha256\_hex}}
\strategy{apikey}{Hashed API keys}{A hotel key card. The front desk does not keep a copy of your card, only a record of which card number opens which room. If someone photographs the desk's records, they still cannot open the door.}
{scripts, integrations and service accounts need a long-lived credential that a person can revoke.}
{the client is a browser (use the session) or can perform an OAuth flow (use short-lived tokens).}
{\code{Crypt::URandom} or \code{/dev/urandom} for 32 bytes, \code{Digest::SHA::sha256_hex} for storage, \code{HSET apikey:<hash>} with \code{EXPIRE}; example 30.}
The server generates the key from a CSPRNG, shows it once, and stores only its SHA-256 (Figure~\ref{fig:apikeyflow}). On every request it hashes the presented key and looks the hash up. Because the lookup is by hash, there is no stored secret to compare and no timing side channel to defend.
\fig{apikeyflow}{API key issue and use. The plaintext key exists only in the response that creates it and in the client's configuration.}
\begin{itemize}
\item \textbf{No salt, no slow hash.} Argon2 and bcrypt exist to protect low-entropy secrets such as passwords from guessing. A 256-bit random key cannot be guessed, so a single SHA-256 is correct, and it keeps authentication to one cheap lookup.
\item \textbf{Recognizable prefix.} \code{slk_} followed by 64 hex digits lets log filters and secret scanners spot a leaked key, and tells support staff which system a key belongs to.
\item \textbf{Shown once, never cached.} The creating response carries \code{Cache-Control: no-store}. Lost keys are replaced, never recovered.
\item \textbf{Revocation and rotation.} \code{user-apikeys:<user>} lists a user's key hashes, so ``revoke all'' is one loop. Allow two live keys per integration so the old one can be retired after the new one is deployed. Expire unused keys.
\item \textbf{Limits per key.} The fixed-window script from the login limiter (Strategy~\ref{str:ratelimit}) runs per key hash. Responses advertise the quota with \code{RateLimit-Policy} and \code{RateLimit} headers in the syntax of the current IETF draft [54], \code{"perkey";q=10;w=60}, plus \code{Retry-After} on 429.
\end{itemize}

% ---------------------------------------------------------------- JWT
\subsection*{Access tokens: short-lived JWTs, verified strictly}
\index{API!Crypt::JWT}
\strategy{jwtrevocation}{Short-lived token with a revocation list}{A day pass at a theme park. Staff at each ride only check the stamp and the date, which is fast. If a pass is reported stolen, its number goes on a short list at the gates until the day is over; after that the pass is useless anyway.}
{many requests must be authorized cheaply, ideally without a lookup, and you still need logout to take effect at once.}
{tokens must live for hours or days (revocation lists grow and stolen tokens stay useful) or every request needs fresh server-side state anyway.}
{\code{Crypt::JWT} \code{encode_jwt}/\code{decode_jwt} with \code{accepted_alg}, \code{verify_exp => 1}, \code{verify_iss}, \code{verify_aud}; \code{SET jwt:deny:<jti> 1 EX <remaining>}; example 30.}
A JWT access token is a signed claim set that a server can verify without asking anyone. That is its advantage and its danger: a valid signature says nothing about whether the token was meant for \emph{this} API, as \emph{this} kind of token, or whether it is still current. Verify all of it.
\CodeFile{perl}{snippets/api-credentials-02.pl}{Access-token verification in \texttt{SessionLab::Tokens}}
\warn{In \code{Crypt::JWT} [52] the default for \code{verify_exp} checks \code{exp} only \emph{if the claim is present}. A token minted without \code{exp}, by a bug in another service or by anyone who holds the key, is accepted forever. Pass \code{verify_exp => 1}. The library does not check the \code{typ} header either, so the explicit access-token type of RFC~9068 (\code{at+jwt}) [50] is checked by hand. Without that check, any other JWT signed with the same key, such as an ID token, would be accepted as an access token.}
Revocation costs two Redis reads: \code{EXISTS jwt:deny:<jti>}, set at logout, and \code{EXISTS rtfam:{fam}}, which disappears when the token's refresh family is revoked. The denylist entry expires when the token would have expired anyway, so the list never grows beyond five minutes of logouts. Drop the family check if you can tolerate up to five minutes of access after revocation; keep it if not.
\note{The laboratory signs with HS256, a shared secret, because one application both issues and verifies. When other services verify your tokens, switch to an asymmetric algorithm (ES256, EdDSA or RS256) with a \code{kid} header and a published JWKS, so verifiers never hold a key that can mint tokens. RFC~8725, the JWT best current practice [49], covers algorithm pinning and explicit typing; an update to it is in progress at the IETF.}

% ---------------------------------------------------------------- refresh
\subsection*{Refresh-token rotation and the concurrency trap}
\index{Refresh-token rotation}
\strategy{refreshrotation}{Refresh-token rotation with reuse detection}{A cinema ticket that is torn when you go in. If someone shows up later with the other half of a torn ticket, the cinema knows a copy exists and cancels every ticket bought on that booking, including yours, until you prove who you are again.}
{long-lived sessions for API clients that cannot keep a secret perfectly: mobile apps, single-page applications, desktop tools.}
{the client can use sender-constrained tokens (mutual TLS or DPoP, RFC~9449 [57]), which make a stolen token useless by itself; or the client is a confidential server with a secure store.}
{one Lua script per refresh; successor = \code{HMAC(secret, old token)}; family hash tag \code{{fam}}; grace window for duplicates; example 30.}
RFC~9700, the OAuth~2.0 security best current practice published in January 2025 [48], requires refresh tokens for public clients to be either sender-constrained or rotated. Rotation means every refresh returns a new refresh token and invalidates the old one. If the old one is ever presented again, someone has a copy: the server revokes the entire chain, called the \emph{token family}, and the user signs in again.

The concurrency trap is that legitimate clients present the same refresh token more than once. Two browser tabs notice an expired access token at the same moment; a mobile app retries after a timeout whose response was lost. Figure~\ref{fig:refreshrace} shows what a read-then-write implementation does with three simultaneous refreshes: every request reads ``active'' before any of them writes, and each mints its own successor. The family now has three live branches, and reuse detection can no longer tell a thief's copy from a second tab. A strict implementation that revokes on the second presentation has the opposite bug: it logs out every user who has two tabs open.
\fig{refreshrace}{Six simultaneous refreshes in the laboratory: the naive version left three live refresh tokens in one family; the atomic version left exactly one, and every caller received it.}
The fix has three parts:
\begin{enumerate}
\item \textbf{One atomic decision.} A Lua script checks the family, reads the token's state, and either rotates it, replays the recorded successor, or revokes the family, all in one step. Redis runs scripts one at a time, so duplicate presentations are serialized exactly as the single-command patterns in Chapter~\ref{ch:redis} are.
\item \textbf{A deterministic successor.} The successor is \code{HMAC(secret, presented token)}, so every duplicate can be handed the same token without the server ever storing a plaintext refresh token. Only the hashes live in Redis.
\item \textbf{A short grace window.} A duplicate presentation within the window receives the same successor (\emph{replay}); after the window, it is treated as theft (\emph{reuse}). The script reads Redis \code{TIME} so that application clocks never matter.
\end{enumerate}
The Perl controller computes the successor and hashes, then passes three keys and three arguments to its Redis wrapper. This application-side excerpt is from the complete controller in Section~\ref{sec:catalyst-lab}.
\CodeFile{perl}{snippets/api-credentials-04.pl}{Perl controller excerpt -- submit the atomic operation}
The next listing runs inside Redis, not in the Catalyst worker. The numbered keys and arguments correspond directly to the invocation above. Its return value selects the subsequent Perl response path.
\CodeFile{lua}{snippets/api-credentials-03.lua}{Redis operation -- refresh-token rotation}
\fig{refreshstates}{The life of one refresh token. Every token in a family shares the hash tag, so the script's keys live in one Redis Cluster slot.}
The refresh token itself carries its family ID (\code{<fam>.<64 hex digits>}), so the server can build the \code{{fam}}-tagged keys before any lookup. Deleting \code{rtfam:{fam}} is the single switch that revokes everything: the next refresh returns \code{revoked}, and the family check makes outstanding access tokens fail at once.
\warn{The grace window is a deliberate trade-off. Within it, an attacker who replays a stolen token receives the same successor as the real client, so detection waits until one of them rotates again. Keep the window short; identity providers expose it as a setting (Okta's default is 30 seconds, configurable from 0 to 60; Amazon Cognito allows 0 to 60 seconds; Auth0 calls it the reuse interval) [58--60]. Never extend it to hide client bugs.}
\tip{Make clients refresh single-flight. In an async Perl client, guard the refresh call with a \code{Future::Mutex} (Section~\ref{sec:async}) and have every waiting request reuse the result; in a browser, share one in-flight promise across tabs with a \code{BroadcastChannel} or a Web Lock. The server's grace window then becomes a safety net, not the main mechanism.}

% ---------------------------------------------------------------- logout
\tested{Example 30 exercised 33 checks under four Starman workers, three times. Only API-key hashes reached Redis; altered keys and invalid token claims were rejected; rate limits returned 429. Atomic refresh produced one successor for six concurrent requests, whereas the naive path left three live tokens.}

\subsection*{Logout and revocation}
Logout must stop three things: the current access token, every future refresh, and any access token already handed to another tab. The laboratory does all three with two commands: \code{SET jwt:deny:<jti> 1 EX <seconds until exp>} and \code{DEL rtfam:{fam}}. A password change or an administrator's ``sign out everywhere'' deletes every family listed for the user, just as \code{logout_everywhere} deletes every session (Section~\ref{sec:catalyst-lab}). Because both keys expire on their own, revocation leaves no residue to clean up.

% ---------------------------------------------------------------- signed
\subsection*{Signed requests and replay protection}
\index{Nonce}\index{Replay protection}
\strategy{signedrequest}{Signed request with a single-use nonce}{A cheque with a serial number. The bank checks the signature first, then the date, and finally that it has never cashed that serial number before. A photocopy of a cashed cheque is worthless.}
{a partner calls you over networks or proxies you do not control, the body must not be altered, and a captured request must not be replayable.}
{TLS between two systems you operate already provides integrity; mutual TLS or RFC~9421 HTTP Message Signatures are available and agreed with the partner.}
{\code{Digest::SHA::hmac_sha256_hex} over a canonical string; a constant-time comparison; \code{SET nonce:<kid>:<nonce> 1 NX EX 600}; example 30.}
The client signs the method, path, timestamp, nonce and raw body with a per-key signing secret. The server checks them in a strict order (Figure~\ref{fig:signedverify}): first everything that needs no writes, then the signature, and only then records the nonce with \code{SET NX}. A replay inside the timestamp window finds the nonce and gets 409; a replay outside it fails the timestamp check. The nonce TTL is twice the window so that a nonce cannot expire while its timestamp is still acceptable.
\fig{signedverify}{Verification order for a signed request. The first write happens only after the signature has been verified.}
\warn{Record the nonce \emph{after} verifying the signature, never before. Otherwise anyone who can see or guess a nonce can send a forged request carrying it, and the legitimate request that follows is rejected as a replay. Compute the signature over the raw request bytes, before any JSON parsing or normalization, and compare signatures in constant time.}
\note{For new partner APIs prefer a standard: RFC~9421 HTTP Message Signatures [56] defines the canonical string, the covered components and the header syntax. The ordering and nonce rules above apply unchanged.}

% ---------------------------------------------------------------- idempotency
\tested{Reuse outside the 1.5-second grace window revoked the family. Logout denied the token for its remaining lifetime. Signed-request replay (409), stale timestamps (401) and forged signatures (401) were rejected; an invalid signature did not consume the nonce.}

\subsection*{Idempotency for API clients}
API clients retry far more often than people double-click, so the idempotency design in Section~\ref{str:idempotency} matters even more here. The OpenAPI document declares the \code{Idempotency-Key} header as a reusable parameter on every unsafe operation, so client generators send it automatically.
\version{The IETF \code{Idempotency-Key} header draft (draft-07) [55] expired in April 2026 without becoming an RFC. The header remains a widely used convention, popularized by payment APIs. Document the exact semantics, such as key format, retention period and the response to a mismatched body, in the API description rather than relying on the draft. The RateLimit header fields draft was revised in May 2026 (draft-11) and is still in progress.}

% ---------------------------------------------------------------- keys
\subsection*{Key layout and access control}
\begin{table}[!htbp]\centering\small
\caption{Redis keys for API credentials. Braces mark Redis Cluster hash tags.}\label{tab:apikeys}
\begin{tabular}{L{.29\linewidth}L{.13\linewidth}L{.17\linewidth}L{.28\linewidth}}\toprule
Key & Type & TTL & Contents \\ \midrule
\code{apikey:<sha256>} & hash & 90 days & user, created \\
\code{user-apikeys:<user>} & set & none & key hashes, for listing and revocation \\
\code{rtfam:{fam}} & hash & 30 days, sliding & user, current token hash; deleting it revokes the family \\
\code{rt:{fam}:<sha256>} & hash & 30 days & state, used\_at, successor, user \\
\code{jwt:deny:<jti>} & string & until token exp & present while a logged-out token is still unexpired \\
\code{nonce:<kid>:<nonce>} & string & 600 s & single-use marker \\
\code{ratelimit:api:<kid>} & string & 60 s & request count in the current window \\
\bottomrule\end{tabular}
\end{table}
All of these keys belong on the \code{noeviction} instance alongside sessions: evicting \code{rtfam} logs users out, and evicting a denylist or nonce entry silently re-enables a revoked token or a replay. The ACL in Section~\ref{sec:redis-security} extends naturally to the new prefixes.

